\section{Chapter~\ref{sec:serocomputer}. \nt{SERO} neurons}

The properties of \nt{SERO} neurons themselves (rhythm, autoinhibition, sign set by the receptor, subsystems) are measured directly; the chapter's computations (regulator, filter, logic) follow from its equations; $\tau_c$, the diffusion coefficient of serotonin and the number of release sites are estimates; the loop's role as a level regulator in the brain is a hypothesis (in the raphe nucleus autoinhibition is point to point and causes only brief pauses), and the role of \nt{SERO} as the horizon is a hypothesis with behavioral experiments in its favor.

{\footnotesize
\setlength{\tabcolsep}{3pt}\renewcommand{\arraystretch}{1.15}
\begin{longtable}{>{\raggedright}p{0.5cm}>{\raggedright}p{4.0cm}>{\raggedright}p{5.6cm}>{\raggedright}p{2.3cm}>{\raggedright\arraybackslash}p{1.7cm}}
\toprule
No. & Claim & Experiment and what was measured & Source & Status \\
\midrule\endfirsthead
\toprule
No. & Claim & Experiment and what was measured & Source & Status \\
\midrule\endhead
1 & The sign of serotonin's action is chosen by the recipient's receptor (Proposition~\ref{prop:receptor}) & Serotonin receptors are divided into families by pharmacology and mechanism: 5-HT$_{1A}$ opens potassium channels through a G protein and inhibits the cell; 5-HT$_{2A}$ and the ionotropic 5-HT$_3$ excite. One transmitter gives opposite responses in different cells. & \cite{BarnesSharp1999} & measured \\
2 & In the waking state a \nt{SERO} neuron steadily fires a few spikes per second & Recordings of dorsal raphe neurons in freely moving cats: slow rhythmic firing of $2.8\pm0.2$ spikes/s in quiet waking; in slow-wave sleep the rate falls by $34$--$68\,\%$, in REM sleep by $98\,\%$. In anesthetized rats, $1.7\pm0.5$ spikes/s, very regular. & \cite{TrulsonJacobs1979, GagnonParent2014, JacobsAzmitia1992} & measured \\
3 & A \nt{SERO} neuron is a pacemaker: it needs no push for each spike, but it needs a steady background drive (in a slice, noradrenaline through $\alpha_1$ receptors; in the organism, from \nt{NORA}) & In a brain slice the cells fire slowly and steadily, with an afterhyperpolarization of $200$--$800\ms$ after each spike. Fewer cells fire spontaneously in a slice than in the organism: the steady rhythm needs a tonic noradrenaline input through $\alpha_1$ receptors, and blocking them abolishes it. & \cite{VandermaelenAghajanian1983} & measured \\
4 & Almost all receptors are metabotropic; the response is slow, rising and decaying within about a second & All receptor families except 5-HT$_3$ are coupled to G proteins. In a slice, evoked serotonin release produces an inhibitory current through 5-HT$_{1A}$ that rises and decays within a second. & \cite{BarnesSharp1999, CourtneyFord2016} & measured \\
5 & In target regions serotonin is released into the volume, not only into the cleft; in the raphe nucleus itself, inhibition of neighbors is point to point & Fast-scan cyclic voltammetry in slices: release per spike is the same in a region with synapses and in the raphe nucleus, where there are almost no synapses; it does not depend on blocking uptake or receptors; there are far fewer molecules per terminal than transporters and receptors, and the extracellular concentration matches the affinity of the main receptor. In the raphe nucleus, inhibition through 5-HT$_{1A}$ is point to point: the transporter keeps serotonin from accumulating outside the cleft. & \cite{Bunin1998, CourtneyFord2016} & measured \\
6 & The concentration in~\eqref{eq:seroneuron} decays because of reuptake; $\tau_c$ of order a second is an estimate & Voltammetry in the living brain: extracellular serotonin is limited primarily by reuptake and degradation, not by synthesis. The cited papers contain no direct measurement of $\tau_c$; its order of magnitude is the chapter's estimate. & \cite{Hashemi2012serotonin} & measured; $\tau_c$ is an estimate \\
7 & NOT and NOR with a single pacemaker; completeness of \nt{SERO} logic (Proposition~\ref{prop:serocomplete}, Fig.~\ref{fig:serogates}) & Follows from~\eqref{eq:seroneuron}: an inhibitable pacemaker is a NOR gate, and NOR is functionally complete. There is no experiment in which logic has been built from \nt{SERO} neurons; the condition of separate volumes does not hold in the brain. & derived in the chapter & theory \\
8 & Serotonin inhibits the \nt{SERO} neuron itself through receptors on it & In anesthetized rats, 5-HT$_{1A}$ agonists given intravenously and iontophoretically inhibit the firing of dorsal raphe neurons with the same dose--response relation as serotonin itself; 5-HT$_{1B}$ agonists have almost no effect. In a slice, endogenous serotonin from neighboring cells produces a current through 5-HT$_{1A}$ and a brief pause in firing without changing the mean rate. & \cite{Sprouse1987, CourtneyFord2016} & measured \\
9 & In the model the loop through a common volume is a level regulator: variation and the time constant shrink by a factor of $1+G$~\eqref{eq:feedback}; that the loop works this way in the brain is a hypothesis & The classical formula of an amplifier with negative feedback, derived again in the chapter. The loop gain $G$ of the \nt{SERO} system has not been measured. Measured: in a slice, autoinhibition is point to point, causes a brief pause and does not change the mean rate. & \cite{Black1934feedback, CourtneyFord2016} & theory; in the brain, hypothesis \\
10 & The concentration is a moving average of the rate, a low-pass filter~\eqref{eq:lowpass}, Fig.~\ref{fig:lowpass} & Solution of the linear equation~\eqref{eq:seroneuron}; the figure is computed from this formula with $\tau_c=1\,\text{s}$. There is no separate model in \texttt{models/}. & derived in the chapter & theory \\
11 & The tortuosity of the gaps slows the diffusion of a small molecule by a factor of about $2.6$ & Point-source method: iontophoresis of tetramethylammonium and recording of its concentration with an ion-selective electrode in slices and in the living brain. Tortuosity $\lambda\approx1.6$, so the effective $D$ is $\lambda^2\approx2.6$ times smaller than the free one; the extracellular volume fraction is about $0.2$. The diffusion coefficient of serotonin itself in tissue was not measured in these studies; in the chapter it is an estimate. & \cite{Sykova2008} & measured \\
12 & The length of an independent ``wire'' is $\ell\sim\sqrt{D\tau}\approx20$~\textmu m~\eqref{eq:diffusionlength}; the number of wires is limited by the number of such regions (Proposition~\ref{prop:volumewires}) & The diffusion law with the estimates of $D$ and $\tau$ substituted (rows~6 and~11); Proposition~\ref{prop:volumewires} is a consequence. There is no direct experiment counting the independent channels of volume transmission. & derived in the chapter & theory \\
13 & The output of a single \nt{SERO} neuron is spread over huge parts of the brain; the number of release sites is estimated at $\sim10^5$ & Complete reconstruction of single dorsal raphe neurons in the rat: all ascending axons branch heavily, the total axon length of one cell reaches $18.7$~cm, and the branches of one cell reach the striatum, prefrontal cortex and amygdala. The number of release sites per cell has not been counted directly. & \cite{GagnonParent2014} & measured; $10^5$ is an estimate \\
14 & \nt{SERO} neurons fall into several subsystems with different outputs and responses & Viral-genetic labeling in mice: neurons projecting to the amygdala and to the frontal cortex hardly overlap in their branching, receive different inputs and respond oppositely to an aversive stimulus; activating and silencing each group changes behavior in different ways. & \cite{Ren2018} & measured \\
15 & Patience condition $D<\tau_\gamma\ln(R/r_0)$~\eqref{eq:patience} with exponential discounting; with the measured, hyperbolic discounting, $D<\tau_\gamma(R/r_0-1)$ & Follows from discounting by $e^{-D/\tau_\gamma}$ or $1/(1+D/\tau_\gamma)$; derived in the chapter. Monkeys choosing at delays of seconds: discounting is closer to a hyperbola than to an exponential; the response of \nt{DOPA} neurons to the cue of a delayed reward also falls hyperbolically with delay. & derived in the chapter; \cite{KobayashiSchultz2008} & theory \\
16 & \nt{SERO} sets the horizon $\tau_\gamma$ (Section~\ref{sec:horizon}) & Optogenetic activation of dorsal raphe \nt{SERO} neurons in mice lengthens waiting for a delayed reward and increases persistence in foraging for a reward that has become scarce. In humans, lowering serotonin (a tryptophan-free diet) makes the discounting of delayed reward steeper. How the concentration becomes $\tau_\gamma$ is unknown. & \cite{Doya2002, Miyazaki2014, Lottem2018, Schweighofer2008} & hypothesis \\
\bottomrule
\end{longtable}}
